\documentclass[11pt]{article}

\usepackage[a4paper,margin=25mm]{geometry}
\usepackage{amsmath}
\usepackage{amssymb}

\title{A $3\times3$ Counterexample to Convexity}
\author{}
\date{}

\begin{document}

\maketitle

This example isolates the source of nonconvexity in the physical CML
data-fidelity term. For one link,
\begin{equation}
    \widehat{A}(R)
    = \sum_{p\in\ell} d_p k R_p^\alpha,
    \qquad
    J_{\mathrm{atten}}(R)
    = \frac{\bigl(\widehat{A}(R)-A_{\mathrm{obs}}\bigr)^2}{L}.
\end{equation}
This is the one-link normalization used by the normalized solver.

\section{Toy Link}

The example has the following parameters:
\begin{itemize}
    \item a $3\times3$ pixel grid;
    \item one horizontal link crossing the three pixels in the middle row;
    \item a crossed length of $0.125\,\mathrm{km}$ in each pixel, giving a
          total link length of $L=0.375\,\mathrm{km}$;
    \item a frequency of $38.003\,\mathrm{GHz}$ with horizontal polarization;
    \item the ITU-R P.838-3 coefficients
          $k=0.400171594750$ and $\alpha=0.881535333268$.
\end{itemize}

Consider two strictly positive rain fields, measured in
$\mathrm{mm}/\mathrm{h}$, that are left-right reflections of each other
across the vertical centerline of the grid. In $R^{(L)}$, the heavy-rain
value is in the middle-left pixel. In $R^{(R)}$, the same value is in the
middle-right pixel:
\begin{equation}
    R^{(L)} =
    \begin{bmatrix}
        1 & 1 & 1 \\
        20 & 1 & 1 \\
        1 & 1 & 1
    \end{bmatrix},
    \qquad
    R^{(R)} =
    \begin{bmatrix}
        1 & 1 & 1 \\
        1 & 1 & 20 \\
        1 & 1 & 1
    \end{bmatrix}.
\end{equation}

Either field could be chosen as the ground truth. If $R^{(L)}$ is the ground
truth, it generates the observation below, but $R^{(R)}$ produces exactly the
same observation. The same statement holds with the roles reversed. The
single horizontal link therefore cannot distinguish these two rain fields.
Because the link crosses the reflected pixels over equal distances, their
predicted attenuations are equal:
\begin{equation}
    \widehat{A}\bigl(R^{(L)}\bigr)
    = \widehat{A}\bigl(R^{(R)}\bigr)
    = A_{\mathrm{obs}}
    = 0.801595407171\,\mathrm{dB}.
\end{equation}
It follows that
\begin{equation}
    J_{\mathrm{atten}}\bigl(R^{(L)}\bigr)
    = J_{\mathrm{atten}}\bigl(R^{(R)}\bigr)
    = 0.
\end{equation}

The arithmetic mean of the two fields is
\begin{equation}
    R^{(M)}
    = \frac{R^{(L)}+R^{(R)}}{2}
    =
    \begin{bmatrix}
        1 & 1 & 1 \\
        10.5 & 1 & 10.5 \\
        1 & 1 & 1
    \end{bmatrix}.
\end{equation}

\section{Why the Power Law Is Strictly Concave}

To see why the power law is strictly concave, define $g(r)=r^\alpha$ for
$r>0$. Its second derivative is
\begin{equation}
    g''(r)=\alpha(\alpha-1)r^{\alpha-2}.
\end{equation}
Here, $\alpha>0$, $\alpha-1<0$, and $r^{\alpha-2}>0$. Therefore,
$g''(r)<0$ for every $r>0$, which means that $g$ is strictly concave. In
particular,
\begin{equation}
    \left(\frac{20+1}{2}\right)^\alpha
    > \frac{20^\alpha+1^\alpha}{2}.
\end{equation}
In the arithmetic-mean field, the two reflected rain rates $20$ and $1$
become $10.5$ and $10.5$. The strict-concavity inequality shows that their
combined attenuation contribution increases. Consequently,
\begin{equation}
    \widehat{A}\bigl(R^{(M)}\bigr)
    = 0.845082934781\,\mathrm{dB}
    > A_{\mathrm{obs}},
\end{equation}
and
\begin{equation}
    J_{\mathrm{atten}}\bigl(R^{(M)}\bigr)
    = 0.005043106820
    > 0.
\end{equation}

Therefore,
\begin{equation}
    J_{\mathrm{atten}}\!\left(
        \frac{R^{(L)}+R^{(R)}}{2}
    \right)
    >
    \frac{
        J_{\mathrm{atten}}\bigl(R^{(L)}\bigr)
        + J_{\mathrm{atten}}\bigl(R^{(R)}\bigr)
    }{2}
    = 0.
\end{equation}
This violates the defining inequality for a convex function. Therefore, the
physical attenuation data term is nonconvex.

\end{document}
